\subsection{有理函数的原函数}

公式 (31) 使我们能够求出实变量 x 的任意有理函数 \(r(x)\) （系数为实的或复的）的原函数。我们知道（A.VII.7），这样的函数可以（以唯一的方式）写成有限多项之和，这些项是：

\begin{enumerate}
\def\labelenumi{\alph{enumi})}
\item
  或者形如 \(ax^{p}\) （p 为整数 \(\geqslant 0\) ，a 为复数）；
\item
  或者形如 \(a/(x - b)^{m}\) （m 为整数 \(\geqslant 0\) ，a 与 b 为复数）。
\end{enumerate}

现在，容易求出其中每一项的原函数：

\begin{enumerate}
\def\labelenumi{\alph{enumi})}
\tightlist
\item
  \(ax^p\) 的原函数是 \(a\frac{x^{p + 1}}{p + 1}\) ；
\end{enumerate}

\(b)\) 若 \(m > 1\) ，\(a / (x - b)^{m}\) 的原函数是 \(\frac{a}{(1 - m)(x - b)^{m - 1}}\) ；

\begin{enumerate}
\def\labelenumi{\alph{enumi})}
\setcounter{enumi}{2}
\tightlist
\item
  最后，由公式 (10)（III, p.~93）与 (31)（III, p.~101），\(\frac{a}{x - b}\) 的原函数是 \(a\log |x - b|\) （若 \(b\) 为实数），是 \(a\log (x - b)\) （若 \(b\) 为复数）。在后一情形，若 \(b = p + iq\) ，则还有（III, p.~100，公式 (30)）
\end{enumerate}

\[
\log (x - b) = \log \sqrt {(x - p) ^ {2} + q ^ {q}} + i \operatorname{Arc} \tan \frac {x - p}{q} \pm i \frac {\pi}{2}.
\]

我们把更实用的、明显地求出显式给出的有理函数的原函数的方法的考察，推迟到本丛书专讲数值计算的部分。

下列求原函数的问题都可以归结为求有理函数的原函数：

\(1^{\circ}\) 求形如 \(r(e^{ax})\) 的函数的原函数，其中 \(r\) 是有理函数而 \(a\) 是实数；事实上，通过变量替换 \(u = e^{ax}\) ，可归结为求 \(r(u) / u\) 的原函数；

\(2^{\circ}\) 求形如 \(f(\sin ax, \cos ax)\) 的函数的原函数，其中 \(f\) 是两个变量的有理函数而 \(a\) 是实数；变量替换 \(u = \tan(ax/2)\) 把它归结为求下列函数的原函数

\[
{\frac {2}{1 + u ^ {2}}} f \left({\frac {2 u}{1 + u ^ {2}}}, {\frac {1 - u ^ {2}}{1 + u ^ {2}}}\right).
\]

